% Generated draft: requires mathematical and source review.
\documentclass[11pt]{article}
\usepackage[margin=25mm]{geometry}
\usepackage{fontspec}
\setmainfont{Latin Modern Roman}
\usepackage{amsmath,amssymb,mathtools,unicode-math}
\setmathfont{Latin Modern Math}
\usepackage{enumitem,needspace,xurl,hyperref}
\setlength{\parindent}{0pt}
\setlength{\parskip}{5pt}
\setlength{\emergencystretch}{3em}
\newcommand{\N}{\mathbb{N}}
\newcommand{\Z}{\mathbb{Z}}
\newcommand{\Q}{\mathbb{Q}}
\newcommand{\R}{\mathbb{R}}
\newcommand{\C}{\mathbb{C}}
\newcommand{\F}{\mathbb{F}}
\newcommand{\A}{\mathbb{A}}
\newcommand{\PP}{\mathbb{P}}
\newcommand{\E}{\mathbb{E}}
\newcommand{\eps}{\varepsilon}
\newcommand{\dd}{\,\mathrm{d}}
\begin{document}
\section*{UnsolvedMath ID: 1900009}
\textbf{Source problem code:} catalogue entry
\subsection{Definitions and mathematical statement}
Arnold's question asks for a decision algorithm for realizability of a prescribed type of fundamental domain by a periodic continued fraction, in the source's geometric continued-fraction and sail setting.

To expose the required algorithmic quantifiers, suppose the intended class of types has a specified finite encoding with valid-code set $\mathcal D$. For $T\in\mathcal D$, let $\mathcal R(T)$ mean that there exists a periodic continued fraction in the intended class having a fundamental domain of type $T$, with type equality interpreted in the intended equivalence relation. The requested algorithm would be one fixed procedure $A$ such that for every valid input $T$,
\[
A(T)\text{ halts, and}\qquad
A(T)=\mathrm{YES}\ \Longleftrightarrow\ \mathcal R(T),
\]
with output $\mathrm{NO}$ otherwise. Both realizable and unrealizable valid inputs must terminate. Merely enumerating some realizations, or computing a domain for an already specified continued fraction, is not the requested decision procedure. Behavior on invalid encodings depends on the missing input convention.

The source does not define its continued-fraction class, periodicity, the acting group and its fundamental domains, or what data a ``type'' retains. In particular it is not stated whether a type is purely a finite combinatorial cell structure or also includes lattice lengths, face incidences, identifications, or other arithmetic labels. Nor does it specify whether dimension is fixed or part of the input. The displayed algorithmic formulation is therefore a schema with these unresolved parameters, not a claim that an arbitrary finite complex is the intended input or that one-dimensional periodic continued fractions capture the question.

The supplied moderate-confidence triage describes classical structural links with totally real algebraic data and torus decompositions, and related computations of Klein polyhedra and periodic continued fractions, but reports no general realizability decision procedure. These are attributed contextual and search-status statements, not a verified proof of openness, undecidability, or existence of an algorithm.
\subsection{Short English statement}
Is there a procedure that always decides whether a proposed fundamental-domain type occurs for a periodic geometric continued fraction? The intended type data, equivalence, dimension, and periodicity are absent, so the input and realizability predicate must first be fixed. The supplied triage reports related structural and computational work but no verified general decision algorithm.
\subsection{Sources}
\begin{itemize}
\item Local source registry: \url{data/unsolvedmath/problems.json}, permanent UnsolvedMath ID $1900009$, source code catalogue entry, including the supplied OPEN-TRIAGE summary. Dataset attribution: Ulam AI and UnsolvedMath Contributors; original-source attribution and rights are retained.
\item Karpenkov, \emph{Open problems in geometry of continued fractions} (2017), Problem 9, attributed to V. Arnold: \url{https://arxiv.org/abs/1712.01450}. The supplied metadata reports source-TeX extraction and July 29, 2026 access, and flags current status and rights for review.
\item The supplied summary mentions Jacobi--Perron and Markov-type algorithms, Arnold's problems, and papers by Karpenkov and others on Klein polyhedra and higher-dimensional continued fractions. It supplies no precise additional bibliography or theorem giving a general decision procedure. None of those references was independently inspected for this draft.
\end{itemize}
\end{document}
